\documentclass{article}
\usepackage[utf8]{inputenc}
\usepackage[T1]{fontenc}
\usepackage{geometry}
\usepackage{hyperref}
\usepackage{listings}
\usepackage{xcolor}
\usepackage[normalem]{ulem}

\geometry{a4paper, margin=1in}

\title{Ollama Installation and Maintenance Guide}
\author{(C) Sujeet Banerjee}
\date{}

\begin{document}

\maketitle

\section*{Overview}
This guide installs Ollama and downloads the models used in the bootcamp: \texttt{llama3.2} (LLM) and \texttt{nomic-embed-text} (embedding model).

\subsection*{System Requirements}

\begin{enumerate}
    \item Windows 10/11 or macOS 13+
    \item 8 GB RAM minimum (16 GB recommended)
    \item ~10 GB free disk space
    \item Internet connection (to download models)
    \item \sout{Docker Desktop is NOT required}
    \item \sout{Python is NOT required for Day 1} 
    \textit{(Required for Day2/Day3)}
\end{enumerate}


\section*{Windows 10/11}
\begin{enumerate}
    \item Visit \url{https://ollama.com/download/windows} and download the Windows installer.
    \item Run the installer.
    \item Open PowerShell or Command Prompt.
    \item Verify installation:
    \begin{verbatim}
ollama --version
    \end{verbatim}
\end{enumerate}

If successful, download the models:
\begin{verbatim}
ollama pull llama3.2
ollama pull nomic-embed-text
\end{verbatim}

Run the LLM:
\begin{verbatim}
ollama run llama3.2
\end{verbatim}

\noindent The Ollama server normally listens on \url{http://localhost:11434}.

\section*{macOS}
\begin{enumerate}
    \item Visit \url{https://ollama.com/download/mac} or install via Homebrew:
    \begin{verbatim}
brew install ollama
    \end{verbatim}
    \item Start Ollama (if needed):
    \begin{verbatim}
ollama serve
    \end{verbatim}
    \item Verify installation:
    \begin{verbatim}
ollama --version
    \end{verbatim}
\end{enumerate}

Download the models:
\begin{verbatim}
ollama pull llama3.2
ollama pull nomic-embed-text
\end{verbatim}

Run the LLM:
\begin{verbatim}
ollama run llama3.2
\end{verbatim}

\section*{Useful Commands}
\begin{verbatim}
ollama list              # List downloaded models
ollama rm llama3.2       # Remove a model
ollama show llama3.2     # Model information
ollama serve             # Start server manually
\end{verbatim}

\section*{Change LLM Hyperparameters}
From within Ollama CLI or Terminal:
\begin{verbatim}
  /?                               List available LLM commands
  /set parameter ...               Set a (hyper)parameter
  /set verbose                     Show LLM stats
  /set quiet                       Disable LLM stats
  /set think <high|medium|low>     Enable thinking
  /set nothink                     Disable thinking
  /set parameter num_ctx <int>          Set the context size
  /set parameter temperature <float>    Set creativity level
\end{verbatim}

\section*{Troubleshooting}
\begin{itemize}
    \item Do not use Office-Laptop (certain ports might be blocked due to office compliance)
    \item If 'ollama' is not recognized, restart the terminal after installation.
    \item Firewall software may block \texttt{localhost:11434}.
    \item Verify the server by opening \url{http://localhost:11434} in a browser; a simple response indicates the service is running.
\end{itemize}

\section*{Ollama Maintenance: Cleaning Up Local Logs, Memory, and Storage}

To ensure Ollama runs smoothly without consuming unnecessary system resources, you can periodically clean up stale data, memory, and storage.

\subsection*{1. Freeing Up Memory (RAM/VRAM)}
When a model is running, it stays loaded in memory for a period of time (typically 5 minutes by default). To immediately unload the model and free up RAM/VRAM, run the \texttt{ollama stop} command:
\begin{verbatim}
ollama stop llama3.2
\end{verbatim}
Alternatively, you can restart the Ollama service or app to clear all loaded models from memory.

\subsection*{2. Cleaning Up Model Storage (Disk Space)}
Downloaded models can take up significant disk space. To remove a model you no longer need:
\begin{verbatim}
ollama rm <model_name>
\end{verbatim}
Example: \texttt{ollama rm llama3.2}

\vspace{0.5em}
\noindent\textbf{Manual Cache Deletion:}
If you want to clear the entire model cache (or if \texttt{ollama rm} fails), you can delete the models directory directly:
\begin{itemize}
    \item \textbf{Windows:} Delete the contents of \texttt{C:\textbackslash Users\textbackslash <YourUsername>\textbackslash .ollama\textbackslash models}
    \item \textbf{macOS:} Delete the contents of \texttt{\textasciitilde/.ollama/models}
\end{itemize}

\subsection*{3. Clearing Local Logs}
Ollama generates logs that can grow over time.
\begin{itemize}
    \item \textbf{macOS:} The log files are usually located at \texttt{\textasciitilde/.ollama/logs/}. You can delete old log files in this directory to free up space.
    \item \textbf{Windows:} Logs are typically found in \texttt{\%LOCALAPPDATA\%\textbackslash Ollama\textbackslash logs} or the \texttt{.ollama} directory.
\end{itemize}

\end{document}
